\documentclass[aspectratio=169,10pt,fontset=fandol]{ctexbeamer}
\usepackage{amsmath,amssymb,mathtools,bm,mathrsfs}
\usepackage{graphicx,booktabs,array,tabularx,multirow}
\usepackage{tikz,pgfplots,xparse,xstring,mdframed}
\usetikzlibrary{arrows.meta,calc,angles,positioning,shapes.geometric,graphs}
\pgfplotsset{compat=1.18}
\usetheme{fduclassroom}
\setbeamerfont{normal text}{size=\normalsize}
\setbeamerfont{title}{size=\large}
\setbeamerfont{subtitle}{size=\normalsize}
\setbeamerfont{frametitle}{size=\large}
\setbeamerfont{framesubtitle}{size=\small}
\setbeamerfont{block title}{size=\normalsize}
\setbeamerfont{block body}{size=\normalsize}
\setbeamerfont{section in toc}{size=\normalsize}
\setbeamerfont{section in head/foot}{size=\footnotesize}
\setbeamerfont{author in head/foot}{size=\footnotesize}
\setbeamerfont{title in head/foot}{size=\footnotesize}
\setbeamerfont{date in head/foot}{size=\footnotesize}
\setbeamertemplate{navigation symbols}{}
\setbeamertemplate{itemize item}{\normalsize\textbullet}
\setbeamertemplate{itemize subitem}{\small\textbullet}
\setbeamercolor{structure}{fg=blue!65!black}
\setbeamercolor{block title}{fg=blue!65!black,bg=blue!8}
\setbeamercolor{block body}{fg=black,bg=blue!3}

% ctex fontset=fandol; system fonts are not redistributed.


\title{数列}
\begin{document}
\section{数列}
% FRAME-ID F001
\begin{frame}[t,fragile]{目录}
\normalsize
{\color{blue}\normalsize 1. 数列\par}\vspace{0.3em}
\end{frame}
% FRAME-ID F002
\begin{frame}[t,fragile]{数列}
\normalsize
\begin{block}{定理}
{\normalsize
% SOURCE-BEGIN U0001 0 21
若 $a_n\to a$，则收敛数列有界。
% SOURCE-END U0001
}
\end{block}
\pause
{\normalsize
\textbf{证明：}
% SOURCE-BEGIN U0002 0 34
由定义取 $\varepsilon=1$，尾项有界，有限前项亦有界。
% SOURCE-END U0002
}
\end{frame}
% FRAME-ID F003
\begin{frame}[t,fragile]{数列}
\normalsize
\begin{block}{例题}
{\normalsize
% SOURCE-BEGIN U0003 0 26
计算 $\lim_{n\to\infty}1/n$。
% SOURCE-END U0003
}
\end{block}
\pause
{\normalsize
\textbf{解：}
% SOURCE-BEGIN U0004 0 8
答案是 $0$。
% SOURCE-END U0004
}
\end{frame}
\end{document}
